%%%%%%%%%%%%%%%%%%%%%part.tex%%%%%%%%%%%%%%%%%%%%%%%%%%%%%%%%%%
% 
% sample part title
%
% Use this file as a template for your own input.
%
%%%%%%%%%%%%%%%%%%%%%%%% Springer %%%%%%%%%%%%%%%%%%%%%%%%%%

\begin{partbacktext}
\part{Muscle Function Overview}
\noindent 
Part II on Muscle Function includes four short chapters directly related to how muscles generate tension. Tension generating mechanisms are the primary focus of Chapter \ref{chp:tension} on Tension. To generate active tension a muscle fiber must be excited - Chapter \ref{chp:excitation} on Excitation introduces the general concepts of an excitable membrane and signaling processes and how these relate to excitation - contraction coupling. Muscle function depends on sequences of excitation, across time and space, that determine the rate and quantity of active muscle tension. Chapter \ref{chp:regulation} on Regulation describes the basic mechanisms of regulation of muscle active tension and tone by the central and peripheral nervous system. With all of that in place, we are ready to turn our attention in Chapter \ref{chp:energetics} on Energetics to the transfer of chemical energy as delivered to the muscle fiber as glucose or fat to mechanical energy in the form of tension. 
\end{partbacktext}